\section{Reactions}
A chemical reaction is the process of bonds breaking and forming to rearrange atoms in molecules. There are the same number of atoms on both sides of a chemical reaction because of the conservation of mass. In AP chemistry there are a few new topics in reactions that are new including redox reactions and their corresponding half reactions, net ionic reactions, and the reaction categories. 

Reactants are the starting materials in a chemical reaction, while products are the substances that are formed as a result of the reaction. Some reactions also contain a catalyst, which is a substance that speeds up the reaction without being consumed in the process. This is accomplished by lowering the activation energy of the reaction.

There are many different arrows used to symbolize reactions, including a standard arrow $\rightarrow$ which indicates a reaction that goes to completion, a double arrow $\rightleftharpoons$ which indicates a reaction goes to equilibrium, and a reversible arrow $\leftrightarrows$ which indicates a reaction that can go in either direction. In addition multiple arrows are used to indicate several steps occur $\substack{\quad\quad\rightarrow\\[-1em]\quad \rightarrow \\[-1em] \rightarrow}$ and a single arrow with a line through it $\nrightarrow$ is used to indicate a reaction that does not occur. It is important to not confuse the double arrow, the equilibrium arrow or the resonance arrow $\leftrightarrow$ which indicates a molecule that can be represented by multiple Lewis structures.

\subsection{Physical vs Chemical Changes}
A physical change has no effect on the atoms or molecules of a substance. Instead a physical change only changes the physical properties of a substance, such as its state of matter, shape, or size. Examples of physical changes include melting, boiling, freezing, cutting, and dissolving. A chemical change on the other hand is a change that results in the formation of new substances with different properties. Chemical changes involve the breaking and forming of chemical bonds between atoms. Examples of chemical changes include combustion, rusting, and digestion. Both involve energy exchanges, but chemical changes typically involve larger energy changes than physical changes. Physical changes are easily reversible, while chemical changes are often irreversible\footnote{It is possible to reverse a chemical change, but the energy required to do so is often nearly unachievable in reality.}
\subsection{Reaction Types}
\subsubsection{Clasical Reaction Types}
The classic reaction types are synthesis, decomposition, combustion, single replacement, and double replacement. In a synthesis reaction, two or more reactants combine to form a single product.
\[\ce{A + B -> AB}\]
In a decomposition reaction, a single reactant breaks down into two or more products.
\[\ce{AB -> A + B}\]
In a combustion reaction, a substance reacts with oxygen to produce energy in the form of heat and light.
\[\ce{C_xH_y + O2 -> CO2 + H2O}\]
A single replacement reaction is a reaction where one element replaces another element in a compound.
\[\ce{A + BC -> AC + B}\]
A double replacement reaction is a reaction where the positive and negative ions of two compounds switch places tor form two new compounds.
\[\ce{AB + CD -> AD + CB}\]
\subsubsection{Modern Reaction Types}
The modern reaction types that are more commonly used in AP chemistry are combustion, acid-base, redox, and precipitation reactions. The combustion reaction is the same as the classic combustion reaction. 

Combustion reactions are exothermic and are organic redox reactions. They always have some hydrocarbon and oxygen as reactants and carbon dioxide and water as products.

Acid-base, redox, and precipitation reactions all occur in solution. Since they occur in solution and have many ions that don't interact in the reaction, we use net ionic equations to represent the reaction. In a net ionic equation, we separate all ions that disassociate, writing out the ions as individual reactants or products. For example $\ce{HCl_{(aq)}}$ becomes $\ce{H+_{(aq)} + Cl-_{(aq)}}$. Then we cancel out any spectator ions that appear on both sides of the equation, leaving only the ions that are directly involved in the reaction. If a compound is not dissolved, it is not disassociated and therefore kept together as a single reactant or product.  
\subsubsection{Redox Reactions}\label{sec:OxidationReduction}
The classical single replacement, synthesis, and decomposition reactions become redox reactions.Redox reactions involve the transfer of electrons between reactants. If an element is oxidized, it loses electrons and if it is reduced, it gains electrons. This may seem counter intuitive but this is because the term reduce refers to the oxidation number, which gets lower with more electrons. A mnemonic used to memorize this is OIL RIG, which stands for Oxidation Is Loss, Reduction Is Gain. In a redox reaction, one reactant is oxidized and another reactant is reduced. The oxidation and reduction half reactions can be written separately to show the transfer of electrons more clearly.
In a half reaction the reactants and products are separated into two reactions, one for oxidation and one for reduction. We then ballance both sides of the half reaction by adding electrons to both sides to balance the charge, and then we can combine the two half reactions to get the overall redox reaction. For the reaction between Iron(ii) and Copper(iv) ions, the half reactions would be
\begin{align*}\ce{Fe^2+ -> Fe^3+ + e-} \\
\ce{Cu^4+ + 3e- -> Cu+}\end{align*}
Then we can multiply the first half reaction by 3 to get the same number of electrons in both reactions
\begin{align*}\ce{3Fe^2+ -> 3Fe^3+ + 3e-} \\
\ce{Cu^4+ + 3e- -> Cu+}\end{align*}
Then we can add the two half reactions together to get the overall redox reaction
\[\ce{3Fe^2+ + Cu^4+ -> 3Fe^3+ + Cu+}\]
It is usually easy to figure out what the products of a redox reaction will be using the common oxidation numbers, and from there figure out the half reactions. There are some which are more difficult however including Permanganate and Dichromate redox reactions. These half reactions along with others only occur in acidic or basic solutions, in which case you can balance them using $\ce{H+}$ or $\ce{OH-}$ ions respectively. For example, the half reaction for the reduction of permanganate in acidic solution is
\[\ce{MnO4- + 8H+ + 5e- -> Mn^2+ + 4H2O}\]
A table containing the common redox half reactions can be found below in \ref{table:StandardPotential}
\subsubsection{Precipitation and Acid-Base Reactions}
Double replacement reactions become either acid-base reactions or precipitation reactions. Acid base reactions involve the reaction of an acid and a base to produce water and a salt. 
Precipitation reactions involve the reaction of two aqueous solutions to produce an insoluble solid called a precipitate. The precipitate is a insoluble compound that forms when two aqueous solutions are mixed together. 
\subsection{Solubility Rules}
The factors that affect solubility include the amount of solvent and solute along with pressure, the substances, and the temperature. The most common solvent is water and is a requirement for many chemical reactions. There are a few general rules for weather an ionic compound is soluble in water. These rules are as follows.
\begin{enumerate}
    \item An ionic compound is always soluble if it contains
    \begin{itemize}
        \item A nitrate ion (\ce{NO3-})
        \item An acetate ion (\ce{C2H3O2-})
        \item An alkali metal cation (\ce{Li+}, \ce{Na+}, \ce{K+}, \ce{Rb+}, \ce{Cs+})
        \item An ammonium ion (\ce{NH4+})
        \item A chlorate ion (\ce{ClO3-}, \ce{ClO4-})
    \end{itemize}
    \item An ionic compound is usually soluble if it contains
    \begin{itemize}
        \item A halide ion (\ce{Cl-}, \ce{Br-}, \ce{I-}) except when paired with \ce{Ag+}, \ce{Hg^2+}, \ce{Pb2^2+}, or \ce{Cd^2+}
        \item A sulfate ion (\ce{SO4^2-}) except when paired with \ce{Ca^2+}, \ce{Sr^2+}, \ce{Ba^2+}, \ce{Ra^2+}, \ce{Pb2^2+}, or \ce{Ag+}
    \end{itemize}
    \item An ionic compound is usually insoluble if it contains
    \begin{itemize}
        \item A carbonate ion (\ce{CO3^2-})
        \item A phosphate ion (\ce{PO4^3-})
        \item A oxide or sulfide ion (\ce{O^2-}\ce{S^2-})
        \item A hydroxide ion (\ce{OH-})
    \end{itemize}
    Except when paired with an alkali metal or ammonium ion, in which case they are soluble.
\end{enumerate}

An ionic compound dissociates if it is soluble or if it is a strong acid or base, otherwise it will not dissociate. These compounds that disassociate are called strong electrolytes, while those that do not are called weak electrolytes. This is because disassociated ions allow water to conduct electricity. 
\subsubsection{Strong vs Weak Acids and Bases}
A strong acid is an acid that completely disassociates in water, while a weak acid is an acid that only partially disassociates in water. Similarly, a strong base is a base that completely disassociates in water, while a weak base is a base that only partially disassociates in water. The strength of an acid or base is determined by its ability to donate or accept protons.
The strong acids and bases are as listed below. 
\begin{table}[h]
\label{StrongAcidsBases}
\centering
\begin{tabular}{c c}
\textbf{Strong Acids} & \textbf{Strong Bases} \\
\hline
\ce{HCl} & \ce{LiOH} \\
\ce{HBr} & \ce{NaOH} \\
\ce{HI} & \ce{KOH} \\
\ce{HNO3} & \ce{RbOH} \\
\ce{HClO4} & \ce{CsOH} \\
\ce{HClO3} & \ce{Ca(OH)2} \\
\ce{H2SO4} & \ce{Sr(OH)2} \\
 & \ce{Ba(OH)2} \\
\end{tabular}
\caption{Strong Acids and Bases}
\end{table}
Some common examples of weak acids and bases include acetic acid (\ce{CH3COOH}), hydrofluoric acid (\ce{HF}), carbonic acid (\ce{H2CO3}), Aluminum hydroxide (\ce{Al(OH)3}), and Iron(iii) hydroxide (\ce{Fe(OH)3}). The number of hydrogens that can be donated by an acid or accepted by a base is the proticity of the acid or base. A monoprotic acid can donate one hydrogen, a diprotic acid can donate two hydrogens, and a triprotic acid can donate three hydrogens. The same applies to bases with accepting hydrogens.
\subsubsection{The Forces Behind Solubility}
The solubility rules come from the interactions between the ions and water molecules and is governed by coulomb's law. 
\begin{equation}
    F = k \frac{q_1 q_2}{r^2}\label{eq:coulombs}
\end{equation}
While AP chemistry does not require you to use coulomb's law in calculations you need to understand how it relates to many key ideas including solubility and periodic trends. Coulomb's law states that the force of attraction or repulsion between two charged particles is directly proportional to the product of their charges and inversely proportional to the square of the distance between them. This means that ions with higher charges will have stronger interactions with water molecules, making them more soluble. 
\subsection{Titrations}
A titration is a laboratory technique used to determine the concentration of an unknown solution by reacting it with a solution of known concentration. In AP chemistry there are two main types of titrations, acid-base titrations and redox titrations. Both use an indicator to determine the end point of the titration, and stoichiometry to determine the concentration of the unknown solution. To calculate the concentration of the unknown solution, write out the chemical reaction that is occurring, and use stoichiometry to convert the volume of solution added to find the moles of solute in the unknown solution, then divide by the volume of the unknown solution to find the molarity.